% !TeX root = main.tex
\subsection*{Transformer and main LV room proposal}
The proposal starts with two 1500 kVA transformers serving separate LV sections, providing 3000 kVA in total and a margin of 327.784 kVA. Floors 1--5 are allocated to section A and floors 6--10 to section B, with each section also receiving half of the central plant/public allowance. Under this assumed equal split, each section carries 1336.108 kVA, corresponding to about 2030 A at 380 V. However, full-load N--1 is not provided, so load shedding is required if one unit is lost. Parallel operation through any coupler is prevented.

Room placement follows the lower-left perimeter of the supplied plan, with a 12430 by 5450 mm substation and an adjacent 12430 by 3500 mm LV room, whose depth is assumed. This location allows external plant access and short LV interfaces, while providing separation from the lift/stair/AHU core and independent LV access. Additionally, a 3 m delivery route within 20 m of unloading is reserved, subject to confirmation of the external site conditions.

As shown in the appendix, space is reserved for HV equipment at 3000 by 1000 mm and transformers at 2200 by 1800 mm, together with six LV cabinets. These dimensions are assumed reservations rather than manufacturer footprints. Their clearance envelopes fit without overlap. COP101 v 16 clause 4.3 requires HV clearances of 1000 mm at the rear, 1800 mm at the front and 750 mm at the sides, together with 900 mm at transformer LV terminals and 750 mm elsewhere. Additionally, the proposal includes 2-hour partitions, sealed penetrations and an FFL 150 mm above external level. Meanwhile, the outward-opening 180-degree doors have shared sweeps and therefore require sequential use.

Ventilation is provided by two 600 by 2000 mm louvres, providing 1.20 m$^2$ at an assumed 50\% free-area ratio, above clause 6.3's 1.12 m$^2$ reference. Thus, the certified free-area ratio must be at least 46.67\%. Fan symbols are conditional provisions, not actual heat-loss calculations. Meanwhile, the LV interfaces, including the assumed tails, are 10.03/10.31 m, below the 20 m limit in 5.5.2.2. Actual equipment, structures, openings and ventilation must satisfy these declared conditions.
